\section{Purpose, Scope, and Policy Basis}

\subsection{Purpose}
This framework aims to:
\begin{enumerate}[label=(\alph*)]
  \item empower faculty and staff to use AI effectively and confidently in their professional work;
  \item improve students' learning experience and prepare them for an AI-enabled workplace and society;
  \item protect academic integrity, privacy, and equity through clear, fair, and proportionate rules; and
  \item provide a lightweight, renewable mechanism for governing AI at the University.
\end{enumerate}

\subsection{Scope and applicability}
The framework applies to all \vsu{} students, faculty, researchers, administrative staff, and external collaborators acting for the University, across all campuses and modalities, when they use AI systems for University work. It covers generative AI (text, image, code, audio, video), AI agents and assistants (e.g., chat-based assistants and coding or workflow agents), analytics and prediction tools, and AI features embedded in other software.

Out of scope: purely personal use of AI that does not involve University work, data, or assessment.

\subsection{Policy basis and alignment}
The framework is designed to align with:
\begin{itemize}
  \item \cmo{}, \emph{Policies, Standards, and Guidelines for the Framework on Artificial Intelligence (AI) in Higher Education} \citep{cmo21}, as disseminated to higher education institutions in Eastern Visayas through \rmo{} \citep{rmo318};
  \item Republic Act No.\ 10173 (Data Privacy Act of 2012), its Implementing Rules and Regulations, and National Privacy Commission issuances \citep{ra10173};
  \item existing \vsu{} policies on academic honesty, research ethics, intellectual property, information security, and student discipline.
\end{itemize}

Where this framework conflicts with national law or CHED issuances, the latter prevail. Where it is silent, existing \vsu{} policies apply. The framework is a university-level instrument that implements the Institutional AI Governance Framework contemplated in Section~14 of \cmo{}.

\begin{table}[ht]
\centering
\caption{Alignment map to \cmo{}.}
\label{tab:alignment}
\small
\begin{tabular}{@{}p{4.2cm}p{10cm}@{}}
\toprule
Framework section & Related provisions of \cmo{} \\
\midrule
Principles (\S\ref{sec:principles}) & Sec.~3 (general provisions); Sec.~4 (national guiding principles 1--8). \\
Governance (\S\ref{sec:governance}) & Sec.~14--18 (Institutional AI Governance Framework, high-risk AI, oversight and approval pathways, avoidance of excessive burdens). \\
Faculty enablement (\S\ref{sec:faculty}) & Sec.~7(2) (faculty academic freedom); Sec.~9(1) (AI as an academic tool); Sec.~21 (faculty AI literacy and governance training). \\
Student learning and academic integrity (\S\ref{sec:student}, \S\ref{sec:integrity}) & Principle~2 (syllabi state when AI use is appropriate, restricted, and disclosed); Sec.~8 (student disclosure); Sec.~9(2) (assessment redesign; detection tools not the sole basis for misconduct). \\
Data privacy (\S\ref{sec:privacy}) & Principle~3 (data and AI safeguards; no confidential material in unsecure public platforms); preamble (RA~10173 and NPC issuances). \\
Equity and literacy (\S\ref{sec:equity}) & Principle~1 (access, inclusion, fairness); Principle~2 (questioning and validating AI outputs); Sec.~3(3). \\
Research and authorship (\S\ref{sec:research}) & Sec.~10--12 (AI-assisted research; disclosure in theses and dissertations; authorship attribution and scholarly responsibility). \\
Tool approval (\S\ref{sec:tools}) & Principle~3 (minimal-, moderate-, and high-impact tools); Sec.~15--17 (high-risk systems and approval). \\
Implementation (\S\ref{sec:implementation}) & Sec.~23--25 (phased rollout, reporting, quality assurance); Sec.~30 (twelve-month default transition). \\
\bottomrule
\end{tabular}
\end{table}

\subsection{Authority and interpretation}
The policy owner of this framework is the Technical Working Group (TWG) for AI Policy and Governance, led by Dr.\ Magdalene C.\ Unajan. Interpretive questions go to the TWG (Section~\ref{sec:governance}).

Approval route: the draft is reviewed concurrently by the University Academic Council (UAC) and the University Administrative Council (UADCO); after their revisions and approval, it is endorsed by the Office of the President to the Board of Regents (BOR) for approval.
